\documentclass{article}

\PassOptionsToPackage{hyphens}{url}
\PassOptionsToPackage{numbers,compress}{natbib}
\usepackage[sglblindworkshop]{neurips_2026}
\workshoptitle{AI at Scale for Clinical Impact (ASCI): Cancer Pathology Foundation Models}

\usepackage[utf8]{inputenc}
\usepackage[T1]{fontenc}
\usepackage{courier}
\usepackage{hyperref}
\renewcommand{\topfraction}{0.9}
\renewcommand{\bottomfraction}{0.7}
\renewcommand{\textfraction}{0.1}
\renewcommand{\floatpagefraction}{0.75}
\setcounter{topnumber}{3}
\setcounter{totalnumber}{4}
\hypersetup{hidelinks}
\usepackage{url}
\usepackage{booktabs}
\usepackage{amsmath}
\usepackage{amssymb}
\usepackage{microtype}
\usepackage{xcolor}
\usepackage{graphicx}
\usepackage{flafter}
\usepackage{float}

% Table highlighting: improvements bold, regressions in a muted red. Both survive greyscale
% printing (weight differs), so the mark is not carried by colour alone.
\definecolor{worse}{HTML}{B91C1C}
\newcommand{\gain}[1]{\textbf{#1}}
\newcommand{\loss}[1]{\textcolor{worse}{#1}}

\title{Contrastive fine-tuning against stain and scanner bias improves robustness and downstream performance across pathology foundation models}

\author{%
  \textbf{Ryan Kim}$^{1,3}$ \quad \textbf{Hassan Muhammad}$^{1,2}$ \quad \textbf{Anish Dulal}$^{1,4}$ \\
  \textbf{Paul S. Scotti}$^{1,2}$ \quad \textbf{Tanishq Mathew Abraham}$^{1,2}$ \\[3pt]
  $^1$MedARC \quad $^2$Sophont \\
  $^3$University of Illinois at Urbana-Champaign \quad $^4$University of Oregon \\
  \texttt{ryankim4@illinois.edu}
}

\begin{document}

\maketitle

\begin{abstract}
Histopathology foundation models achieve strong performance across diverse downstream tasks, yet their representations have been shown to encode scanner- and site-specific biases, creating generalization gaps and potentially limiting real-world utility. We introduce \textsc{Spectra} (Scanner and Pathology Embedding Contrastive Tuning for Robust Alignment), a model-agnostic contrastive learning framework which explicitly encourages scanner-robust representations, to optimize performance and generalization on downstream applications. Using parameter-efficient LoRA adapters trained on matched multi-scanner and multi-stain whole-slide images sourced from the PLISM dataset, our method aligns representations of same-location, different-scanner regions while repelling different-location, same-scanner regions. \textsc{Spectra} improved PathoROB robustness index, gene-expression prediction on HEST, slide-level classification AUC on CPTAC, and PLISM retrieval across held-out scanners and stains for every model.  Overall tile-level classification on THUNDER also improved, with exceptions for segmentation and calibration tasks.  We open-source our  training and evaluation codebases and release open model weights for the fine-tuned versions of OpenMidnight and H-optimus-0.
\end{abstract}

\section{Introduction}

A digital pathology-based cancer diagnosis, biomarker prediction, or prognostic assessment should not depend on the laboratory which prepares the reference tissue specimen or the scanner which digitizes it. Pathology foundation models promise to support diverse clinical tasks, yet strong benchmark performance alone does not establish reliability across real-world applications and their related deployment environments \citep{chen2024uni, huang2025flex}. Most leading models are pretrained on proprietary samples drawn from a single institution or a small number of health systems. Although these collections can be large, their limited institutional diversity may expose models to a narrow range of patient populations, laboratory protocols, stains, and scanner hardware. Even The Cancer Genome Atlas (TCGA), one of the largest public digital-pathology repositories, does not eliminate this risk: despite aggregating specimens across more than 140 geographically distributed institutions within a centralized digitization program, each tissue-source-site code denotes a specific institution--cancer-type pair, creating a perfect association between that code and the corresponding disease cohort, making models susceptible to recovering acquisition site from image features \citep{gdc_tss_codes,dehkharghanian2023biased}. Consequently, learned representations can entangle disease morphology with institution-specific acquisition signatures, allowing downstream adaptors to exploit spurious correlations that fail when transferred to cohorts external to training data \citep{huang2025flex}.

Evidence of this vulnerability is increasingly widespread. PathoROB \citep{koemen2026pathorob}, which evaluates 20 foundation models across 34 medical centres, demonstrates persistent sensitivity to quantifiable non-biological variation. Conventional discrimination metrics can also conceal clinically relevant instability. For example, a controlled study of identical slides digitized on different devices found scanner-dependent calibration shifts even when AUC remained largely stable, with robustness not reliably predicted by model size or pretraining scale \citep{thiringer2026scanner}. In real-world practice, these shifts may move patients across diagnostic, prognostic, or biomarker-triage thresholds despite not having relevant tissue biology driving those associations. These findings motivate representation-learning methods which explicitly preserve disease-relevant information, while suppressing technical and non-biological variation, making performance and calibration across clinical environments primary objectives rather than incidental consequences of scale.

Motivated by this domain-wide problem, we introduce \textsc{Spectra}, a low-resource, model-agnostic fine-tuning framework that improves the robustness of pathology foundation-model representations to scanner and staining variation using parameter-efficient LoRA adapters and a label-free contrastive objective over spatially matched multi-scanner, multi-stain tiles. The objective, batch construction, and checkpoint-selection rule are described in \S\ref{sec:methods}.

We fully open-source \textsc{Spectra}, including the complete training and evaluation pipeline, and release fine-tuned H-optimus-0 and OpenMidnight checkpoints intended for direct substitution into existing pipelines.

\section{Related Work}
\label{sec:related_work}

\paragraph{Robustness interventions in computational pathology.}
Existing approaches to addressing model robustness operate at several stages of the computational pathology pipeline. Image preprocessing and augmentation methods reduce colour variation through stain normalization \citep{reinhard2001,macenko2009,vahadane2016structure} or expose models to broader appearance distributions through stain and colour augmentation \citep{tellez2019quantifying,shen2022randstainna}. Although effective against some staining shifts, these methods do not explicitly address scanner-dependent differences in optics, resolution, compression, or image-processing software, and may suppress diagnostically relevant colour information. Training-time approaches use domain-adversarial objectives \citep{ganin2016dann} or information-bottleneck formulations \citep{huang2025flex} to discourage prediction of acquisition domains, but generally require site labels and are optimized for a particular downstream task. Post-hoc methods instead harmonize frozen embeddings through feature alignment or batch correction \citep{nguyen2026fmmap,koemen2026pathorob}; these methods can improve robustness without modifying the encoder but may depend on target-domain statistics and may require adaptation to a new deployment setting. Knowledge distillation offers another route to transferring invariant representations \citep{filiot2025distill}, although it requires a teacher and an additional distillation stage and introduces additional computational cost. Collectively, these approaches mitigate particular forms of domain shift but do not provide a reproducible, task-agnostic mechanism for directly adapting an existing foundation-model encoder to suppress both scanner and staining variation.

\paragraph{Contrastive adaptation using matched scanners.}
To our knowledge, Carloni et al.\ \citep{carloni2025scangen} were the first to introduce contrastive fine-tuning for scanner generalization in pathology-foundation-model pipelines. Their ScanGen objective uses matched specimens digitized on different scanners as positive pairs and different specimens acquired on the same scanner as negative pairs. ScanGen trains a projection network placed between a frozen foundation-model encoder and a multiple-instance learning model, jointly optimizing scanner invariance and supervised EGFR-mutation prediction. The study provided an important proof of concept across five foundation models and a proprietary cohort of 323 lung-cancer specimens digitized on six scanners. However, its evaluation was confined to a single cancer type and biomarker, and robustness was assessed primarily through scanner-dependent variation in EGFR predictions. It did not evaluate whether the adapted representations retain general utility across standardized foundation-model benchmarks, unrelated disease contexts, or a broad range of downstream tasks. Moreover, because adaptation occurs in a task-specific projection network while the encoder remains frozen, the resulting model is not a reusable robust foundation-model backbone. \textsc{Spectra} moves this contrastive adaptation into parameter-efficient LoRA modules within the encoder, incorporates both scanner and staining variation, requires no downstream labels, and evaluates robustness and representation retention across PathoROB, HEST, THUNDER, and a 38-task CPTAC panel.

\paragraph{Task-agnostic foundation-model robustification.}
Recent work by Filiot et al.\ \citep{filiot2026robustifying} reports a model-agnostic fine-tuning procedure applied to ten pathology foundation models. The study evaluates the resulting encoders with PathoROB, HEST, THUNDER, and Patho-Bench and reports improvements in both robustness and aggregate downstream performance. It released two adapted checkpoints: Phaet, derived from Phikon-v2, and Mascaret, derived from Midnight-12k. However, the current preprint does not include a method section or specify the optimization objective, robustness-training data, construction of training pairs, trainable parameters, or optimization hyperparameters. At the time of writing, no training implementation has been released. Consequently, the reported procedure cannot be independently reproduced, audited, or applied to the other eight evaluated backbones. Furthermore, only two adapted checkpoints are available, one of which is Phikon-v2, an already underused model within the larger ecosystem.

\textsc{Spectra} addresses this reproducibility and accessibility gap by providing a fully specified, label-free InfoNCE objective, parameter-efficient LoRA adaptation, and an open training pipeline based on the public PLISM matched-acquisition dataset \citep{ochi2024plism}.

\section{Methods}
\label{sec:methods}

\subsection{Overview}

\textsc{Spectra} is a model-agnostic, parameter-efficient fine-tuning framework which utilizes registered views of the same tissue acquired under different scanner and staining conditions to improve the robustness of a base pathology foundation model. For each model, we freeze the original backbone and train LoRA adapters on the PLISM dataset (\S\ref{sec:data}--\ref{sec:objective}). Standard pathology benchmarking datasets are used exclusively for downstream evaluation and do not influence training; PathoROB alone is used to select the checkpoint (\S\ref{sec:selection}).

\subsection{Training Data}
\label{sec:data}

We use the publicly available PLISM dataset \citep{ochi2024plism}, although \textsc{Spectra} is not specific to this dataset and requires only spatially matched tissue images acquired under different conditions. PLISM was designed to isolate acquisition-related variation using serial sections of a tissue microarray containing 46 human tissue types, stained under 13 distinct H\&E protocols and imaged using seven whole-slide scanners and six smartphone cameras \citep{ochi2024plism}. It comprises two spatially registered subsets: PLISM-wsi, containing 91 whole-slide images representing every combination of the 13 stains and seven scanners, and PLISM-sm, containing corresponding fields of view acquired using both scanners and smartphones. We use only PLISM-wsi in this study.

We define each scanner--stain pair as an acquisition condition
$c=(\text{scanner},\text{stain})$. Rather than extracting tiles directly from the original whole-slide images, we use Owkin's re-registered tile release derived from PLISM-wsi \citep{filiot2025distill}. In this release, each of the 91 whole-slide images is represented by the same ordered set of $N=16{,}278$ spatial locations, extracted as $224\times224$-pixel tiles. We denote the tile from location $i$ under condition $c$ by $x_{c,i}$; consequently, $x_{c,i}$ and $x_{c',i}$ depict corresponding tissue regions under different acquisition conditions. Across scanners at a fixed stain, corresponding tiles originate from the same physical tissue section. Across stains, they originate from adjacent serial sections, which preserve tissue identity and local architecture but do not contain exactly the same cells \citep{ochi2024plism}.

To evaluate generalization beyond the acquisition conditions observed during training, we hold out two scanners (GT450, S210) and three stains (HRH, KR, MY). The remaining five scanners and ten stains define 50 training conditions, while the other 41 conditions contain at least one held-out scanner or stain; the split is fixed across seeds, and tissue locations are shared between training and held-out acquisition conditions (Appendix~\ref{app:evaluation}).

\subsection{Base Models}
We apply \textsc{Spectra} to a diverse set of pathology foundation models to assess the consistency of its improvements across ViT backbones and pretraining histories and to compare robustness and downstream performance before and after fine-tuning. We evaluate Phikon-v2 \citep{phikonv2}, Midnight-12k \citep{midnight}, Virchow2 \citep{virchow2}, H-optimus-0 \citep{hoptimus0}, UNI2-h \citep{uni2h,chen2024uni}, Virchow \citep{virchow}, and OpenMidnight \citep{openmidnight}.

\subsection{Grid-batch construction}
\label{sec:grid}

At each optimizer step we sample $C$ training conditions and $T$ locations, both without replacement, giving a batch of $CT$ images. For an ordered pair $(c_a,c_b)$, $a\neq b$, and location $i$
(Figure~\ref{fig:grid}), the query is $x_{c_a,i}$, the candidate set is the row of $c_b$ alone; the
positive is $x_{c_b,i}$, and the $T-1$ other tiles in that row are negatives. 

Both orderings are included. All candidates in a softmax share one acquisition condition, so acquisition identity alone cannot identify the matching tissue location. The recipe uses $C=2$,
$T=900$, so for each step we have 1,800 images and 899 negatives per query
(alternative allocations of $C$ and $T$ are ablated in \S\ref{sec:ablations}).

\begin{figure}[t]
\centering
\includegraphics[width=\linewidth]{figures/grid_batch.pdf}
\caption{Left: grid batch on PLISM tiles; columns are registered locations, rows are conditions. Right: for the query, the registered tile of another condition is pulled in and the other $T-1$ tiles of that condition are pushed away.}
\label{fig:grid}
\end{figure}


\subsection{CLS and patch-token contrastive objectives}
\label{sec:objective}
Rank-32 LoRA \citep{lora} adapters on the attention and MLP blocks train under $5\%$ of the backbone parameters ($2.7$--$4.4\%$, adapters only). The LoRA-adapted backbone maps a tile $x_{c,i}$ to a CLS token and $P$ patch tokens
$u_1,\dots,u_P$ (register tokens discarded). This gives the two views: $v^{\textsc{cls}}_{c,i}$
is the CLS token, and $v^{\textsc{gem}}_{c,i}$ is a softplus generalized mean \citep{gem} of the
patch tokens with learned exponent $\rho$ (Appendix~\ref{app:pseudo}, Eq.~\eqref{eq:gem}). Each view passes through
its own two-layer MLP head and is $\ell_2$-normalized, giving unit vectors $z^h_{c,i}$.

For an ordered pair of conditions $(c_a,c_b)$, $a\neq b$, query $i$ and candidate $j$, let
$\ell^h_{abij}=\langle z^h_{c_a,i},z^h_{c_b,j}\rangle/\tau$ be their temperature-scaled cosine similarity. Each view $h$ is trained with the same reweighted InfoNCE loss
\citep{infonce},
\begin{equation}
\mathcal{L}_h=-\frac{1}{C(C-1)T}
\sum_{a\neq b}\sum_{i=1}^{T}
\log\frac{\exp(\ell^h_{abii})}{\sum_{j=1}^{T}w^h_{ij}\exp(\ell^h_{abij})},
\end{equation}
where the weights $w^h_{ij}$ apply pseudo-tissue reweighting~\eqref{eq:w}. Each location $i$ carries a fixed
pseudo-group $q_i$: locations are grouped into the eight-connected spatial components of the tile
lattice (putative tissue-microarray cores), which are then merged to 46 groups by clustering
frozen Virchow2 embeddings of one reference condition, so that a pseudo-group approximates one
tissue type (Appendix~\ref{app:pseudo}). A candidate sharing the query's pseudo-group ($q_j=q_i$,
$j\neq i$) is dropped from the GeM softmax and up-weighted by $e^{3}$ in the CLS softmax, while
all other candidates and the positive keep weight $1$ (ablations in \S\ref{sec:ablations}). The two views are weighted equally,
\begin{equation}
\mathcal{L}=0.5\,\mathcal{L}_{\textsc{cls}}+0.5\,\mathcal{L}_{\textsc{gem}}.
\end{equation}


\subsection{Fine-tuning and checkpoint selection}

Only the adapters, the projection heads, and the GeM exponent are trained; the backbone is
otherwise frozen. The heads and GeM exponent are discarded after training. Evaluation uses the
backbone's original embedding and each benchmark's pooling rule. Training runs 500 steps with
AdamW \citep{adamw} (Appendix~\ref{app:implementation}). Adapters are merged into the base
weights.

\paragraph{Checkpoint selection.}\label{sec:selection}
We perform three independent adapter fine-tuning runs per backbone and select one checkpoint per seed by PathoROB RI
(defined in \S\ref{sec:experiments}) with a tolerance rule ($\epsilon=0.007$, Appendix~\ref{app:selection}). All 21 selected checkpoints fall at steps $100$--$200$, within the 200-step linear warmup.


\section{Experiments}
\label{sec:experiments}

We test whether adaptation against scanner and staining variation improves the usefulness of pathology foundation-model representations across downstream tasks. Our evaluation spans robustness to site and acquisition shifts (PathoROB), cross-condition tissue retrieval (PLISM), and downstream benchmarks, including spatial transcriptomics prediction (HEST), CPTAC classification, and general performance (THUNDER). Full evaluation details are provided in Appendix~\ref{app:evaluation}.

\paragraph{PathoROB.}
We evaluate centre-level robustness using PathoROB across TCGA $2\times2$, Camelyon, and Tolkach-ESCA, and report the robustness index (RI) averaged across the three benchmarks. Let $SO$ denote the number of same-class, other-centre $k$-NN pairs and $OS$ the number of other-class, same-centre pairs, pooled across queries. We compute
\begin{equation}
\mathrm{RI} = \frac{SO}{SO+OS}.
\end{equation}
Higher RI indicates representations that preserve class similarity over centre-specific similarity.

\paragraph{Held-out PLISM retrieval.}
We test cross-acquisition alignment on held-out PLISM data. For each query tile, we perform top-1 retrieval of its registered counterpart among 256 candidates. Scanner-pair evaluations vary the scanner while holding stain fixed; stain-pair evaluations vary the stain.

\paragraph{HEST.}
We follow the HEST evaluation protocol \citep{hest}. We fit PCA--Ridge probes for 50 genes across nine tasks and report mean Pearson correlation $r$ across tasks.

\paragraph{CPTAC.}
We evaluate the CPTAC subset of Patho-Bench \citep{pathobench}, reporting macro one-vs-rest AUC across 38 classification tasks.

\paragraph{THUNDER.}
We use the full THUNDER evaluation suite \citep{thunder} across 16 classification and four segmentation datasets. We report $k$-NN and linear-probe macro-F1 performance, few-shot classification, segmentation, calibration via expected calibration error (ECE), and adversarial robustness.

\paragraph{Evaluation protocol.}
For every benchmark, we compare each base model against its fine-tuned counterpart using the same evaluation harness and the PathoROB-selected checkpoint for each adapter seed. Fine-tuned results are reported as mean $\pm$ two sample standard deviations across three adapter seeds, while base-model results are a single evaluation (deterministic except segmentation, Appendix~\ref{app:evaluation}). Per-task HEST and per-dataset THUNDER seed spreads are tabulated in Appendix Tables~\ref{tab:hestranks}, \ref{tab:ranks} and~\ref{tab:ranks-b}.


\section{Results}
\label{sec:results}

Across benchmarks, fine-tuning consistently improves robustness while largely preserving, and in most settings improving, downstream task performance. The clearest effects are on PathoROB and CPTAC, where improvements occur across all backbones. HEST likewise improves on average for every backbone, although sub-tasks have more mixed results. THUNDER's classification and few-shot performance are generally improved, while calibration shows small regressions for several models and segmentation differences are largely within evaluation variability. Figure~\ref{fig:b2t} summarizes these base-to-fine-tuned changes, while Table~\ref{tab:overall} places them in the context of each benchmark's published model field.

\begin{figure}[H]
\centering
\includegraphics[width=\linewidth]{figures/base_to_tuned.pdf}
\caption{Base $\to$ fine-tuned per backbone, ordered by base PathoROB RI. Hollow: base. Filled: fine-tuned mean over three seeds at the selected checkpoint, bar $\pm2$ SD. Numbers are $\Delta$ vs base. Blue improved, red worse, grey within $2$ SD.}
\label{fig:b2t}
\end{figure}

\begin{table}[htbp]
\caption{Leaderboard position, base $\to$ fine-tuned, in each benchmark's published field:
PathoROB by mean RI, HEST by mean Pearson $r$, THUNDER by rank sum over six tasks, including the
adversarial-drop task reported only in the appendix (Appendix Table~\ref{tab:ranks}). Ties share a
position. Fields hold 23, 26 and 32 models. $^\dagger$: absent from the published field and
inserted into it (24 models) rather than substituted (Appendix~\ref{app:evaluation}).}
\label{tab:overall}
\centering
\small
\setlength{\tabcolsep}{4pt}
\renewcommand{\arraystretch}{0.95}
\input{tables/overall_ranks}
\end{table}

\subsection{Cross-centre robustness: PathoROB}

\begin{table}[htbp]
\caption{PathoROB mean RI, RI rank sum and dense-rank position in the 23-model published field,
base $\to$ fine-tuned ($\pm$ two SDs over three seeds).
$^\dagger$: rank estimated by insertion (Appendix~\ref{app:evaluation}).
Last column: held-out PLISM top-1 retrieval (256 candidates, scanner\,/\,stain; two SDs $\leq0.036$); its held-out axis is acquisition, not tissue.}
\label{tab:pathorobranks}
\centering
\small
\setlength{\tabcolsep}{4pt}
\renewcommand{\arraystretch}{0.95}
\input{tables/pathorob_ranks}
\end{table}

Fine-tuning improves cross-centre robustness for every backbone. Mean RI increases in all cases, by $+0.048$ to $+0.366$, each exceeding the fine-tuned $\pm2$\,SD spread (Figure~\ref{fig:b2t}, Table~\ref{tab:pathorobranks}), with the largest per-dataset gains occurring on Camelyon (Appendix Table~\ref{tab:pathorob-sub}). Across models, fine-tuned RI values also tend to converge toward approximately $0.9$, suggesting diminishing gains among backbones that already begin with relatively strong centre-robust representations.

These improvements in RI do not correspond to a substantial loss in predictive structure. PathoROB's prediction-performance submetric, pooled $k$-NN class purity, decreases on TCGA for every backbone by $0.010$--$0.030$, while $k$-NN balanced accuracy averaged over the three datasets changes by only $-0.005$ to $+0.005$ (on TCGA alone, $-0.014$ to $+0.009$; Appendix Table~\ref{tab:pathorob-sub}). Thus, the large changes in robustness are accompanied by comparatively small changes in local class-discriminative performance.

\subsection{Held-out PLISM retrieval}

Top-1 retrieval of the registered counterpart improves for every backbone on both axes
(Table~\ref{tab:pathorobranks}): $0.703$--$0.978$ to $0.948$--$0.997$ across scanners, and
$0.486$--$0.830$ to $0.782$--$0.931$ across stains. Stain remains the harder axis both before
and after (\S\ref{sec:data}).


\subsection{HEST and CPTAC}

\begin{table}[htbp]
\caption{HEST/CPTAC performance: base $\to$ fine-tuned. HEST is mean Pearson $r$ over nine tasks with the
rank sum over those tasks and the positional rank in the 26-model published field; CPTAC is
macro-OvR AUC over 38 Patho-Bench classification tasks. Per-task results are in
Appendix Tables~\ref{tab:hestranks}, \ref{tab:cptac-per-task} respectively to HEST and CPTAC.}
\label{tab:hestcptac}
\centering
\small
\setlength{\tabcolsep}{4pt}
\renewcommand{\arraystretch}{0.95}
\input{tables/hest_cptac}
\end{table}

The robustness gain is accompanied by gains on both downstream benchmarks. Mean HEST Pearson $r$ increases for every backbone, by $+0.002$ to $+0.017$, while CPTAC macro-OvR AUC increases for every backbone, by $+0.009$ to $+0.038$ (Figure~\ref{fig:b2t}, Appendix Table~\ref{tab:hestranks}, Appendix Table~\ref{tab:cptac-per-task}).

The HEST improvements are not uniform across individual gene-prediction tasks: some backbone-task pairs improve while others regress. The mean change nonetheless exceeds the fine-tuned $\pm2$\,SD spread on six of the seven; on Virchow ($+0.002$ against $\pm0.004$) it lies within seed spread. Fine-tuning therefore does not degrade, and on most backbones modestly improves, mean transcriptomic prediction. Against the per-task seed spread (Appendix Table~\ref{tab:hestranks}), 33 of the 43 per-task improvements and 8 of the 20 regressions exceed $\pm2$\,SD. All seven CPTAC gains exceed the fine-tuned $\pm2$\,SD spread; the smallest is Virchow2 ($+0.009$ against $\pm0.002$) and the widest spread is UNI2-h ($+0.023$ against $\pm0.006$).

\subsection{THUNDER}

THUNDER provides a broader and less uniformly positive picture than PathoROB, HEST, or CPTAC; aggregate task means and ranks are in Appendix Table~\ref{tab:ranks}, per-dataset changes in Appendix Figure~\ref{fig:thunderds}. $k$-NN macro-F1 improves on six of seven backbones ($-0.3$ to $+3.5$ points) and linear-probe macro-F1 on all seven ($+0.1$ to $+1.7$). Few-shot performance varies the most, from $-1.0$ to $+22.3$ points: the $+22.3$ is OpenMidnight, whose base few-shot score of $43.6$ ranks 30th in the 32-model field, and the only decrease, UNI2-h ($-1.0$), equals its $\pm1.0$ seed spread (per-dataset spreads in Appendix Table~\ref{tab:ranks}). The rank sum improves for every backbone except UNI2-h, whose sum worsens from 17 to 23 while it keeps first position. Segmentation differences are comparatively small and lie within the repeat variability of the THUNDER harness (Appendix~\ref{app:evaluation}); because segmentation evaluation is non-deterministic, our measured base values can themselves differ slightly from those reported by THUNDER.

Calibration is the main metric that trends negatively. Mean ECE increases for four of seven backbones, beyond the $\pm2$\,SD seed spread on Midnight-12k, Virchow2 and UNI2-h; the degradation is concentrated on BRACS, which worsens on six of seven and beyond $\pm2$\,SD on Midnight-12k, Virchow2 and Virchow (Appendix Figure~\ref{fig:thunderds}, Appendix Table~\ref{tab:ranks-b}; Calibration paragraph of Appendix~\ref{app:evaluation}). ECE in this evaluation is estimated from relatively few populated calibration bins, making small changes sensitive to the distribution of predictions across bins.

The most consistent classification regression occurs on TCGA-uniform, which decreases in 19 of the 21 cells formed by the seven backbones and the $k$-NN, linear-probe, and few-shot tasks, by $4.6$ points on average over those 19 cells. Together with the TCGA purity decrease in PathoROB, this is consistent with the objective doing what it is designed to do (\S\ref{sec:discussion}).


\section{Ablations}
\label{sec:ablations}

We ablate the three recipe components---pseudo-tissue reweighting, batch composition, and adapter
rank---with one seed per arm: reweighting on Phikon-v2, Midnight-12k, Virchow2, H-optimus-0, and
UNI2-h, the other two on Phikon-v2 (rank also on Virchow2); values are in
Appendix~\ref{app:ablations}, Appendix Table~\ref{tab:ablations}. Reweighting
(Appendix~\ref{app:pseudo}) versus none, at the reweighted arm's selected checkpoint, moves RI by
$+0.009$ on Phikon-v2 and by less than the $\epsilon=0.007$ selection tolerance on the other four,
so we do not claim it as a general mechanism; it stays in the recipe because it helps the weakest
base and is RI-neutral elsewhere. Batch composition
(\S\ref{sec:grid}; $CT=3600$, compared at step 750) favours few conditions and many locations: RI at step 750
rises monotonically from $0.802$ to $0.824$ as $C$ falls from 24 to 2. Adapter rank
(\S\ref{sec:objective}) shows no consistent monotone trend over $r\in\{8,\ldots,128\}$; $r=128$ is lowest on
Virchow2. No sweep moves RI by more than $0.03$, well below the recipe's own
$0.048$--$0.366$ gain, consistent with \S\ref{sec:discussion}: the gain comes from a small, early
adapter update rather than from any one component.

\section{Discussion}
\label{sec:discussion}

Taken together, the benchmarks give a consistent picture: cross-centre robustness and held-out
cross-acquisition retrieval improve on every backbone, both downstream benchmarks move with them,
and THUNDER, the only mixed benchmark, mixes by task rather than by backbone.

Two properties of the procedure are worth noting. First, every selected checkpoint lands at steps
$100$--$200$ (\S\ref{sec:selection}). The robustness gain therefore comes from a small, early
update to the adapters, made while the learning rate is still ramping to its peak; running the
full 500-step schedule is not what produces it. Second, because PathoROB is the selection criterion, its improvement is not an independent
finding, whereas HEST, CPTAC, and THUNDER are an unbiased readout of
what the RI-selected checkpoints do to downstream utility.

The regressions on TCGA-derived tasks (\S\ref{sec:results}) are likely since TCGA tissue-source-site codes are perfectly associated with cancer type, so a representation that stops encoding acquisition site also loses a shortcut that those tasks reward. Whether that trade is acceptable depends on the deployment setting, and it is one reason we report the per-task tables rather than only the aggregate ranks.

\section{Limitations}

\textsc{Spectra} is trained using a single tissue-microarray dataset acquired at one institution. Although PLISM provides unusually well-controlled scanner and staining variation, its specific contribution to \textsc{Spectra}'s performance remains undetermined because we do not compare training on alternative matched-acquisition datasets. Consequently, we cannot establish whether PLISM is optimal or whether larger, more biologically diverse, or multi-institutional datasets would produce greater robustness and downstream transfer. The observed improvements may also depend on PLISM-specific tissue composition, staining protocols, and scanner coverage.

Registration across stains introduces an additional limitation: stain-matched pairs are drawn from adjacent serial sections (\S\ref{sec:data}), which may exhibit genuine local biological differences. Treating these regions as positive pairs therefore introduces correspondence noise and may encourage the model to suppress biologically meaningful variation. More precise cross-section registration, explicit uncertainty modeling, or objectives that distinguish exact from approximate correspondences could reduce this risk.

Three further caveats apply to the evidence, all of which measure representation consistency across centres and probe performance, not diagnostic accuracy. All ablations are single seed (\S\ref{sec:ablations}), and several of the reported differences are smaller than the selection tolerance $\epsilon$ (\S\ref{sec:selection}), so the pseudo-tissue reweighting, batch-composition, and adapter-rank trends should be read as indicative rather than established; LoRA versus full fine-tuning, one versus two views, temperature, and the number of training conditions are untested. PathoROB is used for checkpoint selection, so its reported gains are optimistically biased relative to the other benchmarks. Finally, the improvements are not uniform: $k$-NN class purity on TCGA, TCGA-uniform classification, and calibration on several backbones regress, and users should check the per-task tables for the tasks they care about before substituting a fine-tuned checkpoint.

\section{Conclusion}

Registered-view LoRA contrastive adaptation on PLISM raises PathoROB RI on seven pathology
encoders, with HEST and CPTAC gains and task-dependent THUNDER changes at PathoROB-selected
checkpoints. We release the training and evaluation pipeline and the fine-tuned H-optimus-0 and
OpenMidnight weights so that the procedure can be reproduced, audited, and applied to other
backbones.

\bibliographystyle{plainnat}
\bibliography{references}

% Appendix for neurips_2026.tex. Split out for clarity; \input from the main file
% immediately after the bibliography. Do not add \documentclass or \begin{document}
% here -- this file is included, not compiled on its own.
\appendix

\section{Method details}
\label{app:pseudo}

\paragraph{GeM view.} With $P$ patch tokens $u_1,\dots,u_P$ and learned exponent
$\rho\in[1,8]$, the \textsc{gem} view is
\begin{equation}
v^{\textsc{gem}}_{c,i}
=\Big(\frac{1}{P}\sum_{p=1}^{P}\operatorname{softplus}(u_p)^{\rho}\Big)^{1/\rho},
\label{eq:gem}
\end{equation}
softplus and both powers in~\eqref{eq:gem} acting elementwise.

\paragraph{Reweighting.} Each location $i$ carries a fixed pseudo-group $q_i$, meant to represent
the same tissue type. The softmax weight of candidate $j$ against query $i$ for head $h$ is
\begin{equation}
w^h_{ij}=\begin{cases}
\omega_h,&q_i=q_j\ \text{and}\ i\neq j,\\
1,&\text{otherwise,}
\end{cases}
\label{eq:w}
\end{equation}
with $\omega_{\textsc{cls}}=e^{3}\approx20$ and $\omega_{\textsc{gem}}=0$. The positive has
$w^h_{ii}=1$ and is never reweighted; $\omega_h=1$ in~\eqref{eq:w} recovers standard InfoNCE. This up-weights same-pseudo-group negatives in the CLS loss and excludes them from the GeM loss.

\paragraph{Map construction.} We construct the same pseudo-groups for all models before training. First, the row and column stored in every tile
key are placed on a binary occupancy lattice. 
Directionally connected components along all eight directions are merged to create contiguous spatial regions. 
A component smaller than one quarter of the
median component size is assigned to the nearest larger component by centroid distance. This
geometry-only stage results in 49 groupings, which is 3 more than the 46 tissue types.

We then embed every tile from one reference condition (GMH stain, S60 scanner) with frozen
Virchow2, concatenating its CLS and mean-patch embeddings. Embeddings are averaged and
$\ell_2$-normalized within each spatial component, and cosine distances are computed between
the 49 component means. Average-linkage hierarchical clustering is cut at 46 groups, using
PLISM's documented category count as the intended grouping and thereby forcing three
merges. This does not guarantee that the recovered components correspond one-to-one
with the 46 categories. Every coordinate inherits its component's cluster.

\section{Checkpoint selection}
\label{app:selection}

Let $R_t$ be PathoROB average RI at checkpoint index $t$ (evaluated every 50 optimizer steps) and $B_t=\max_{s\le t}R_s$. Per seed,
\begin{align}
t_{\mathrm{stop}} &= \min\{\,t\ge 2 : R_t-B_{t-1}\le\epsilon\,\}, \\
t^\star &= \min\{\,t\le t_{\mathrm{stop}} : R_t\ge B_{t_{\mathrm{stop}}}-\epsilon\,\},
\end{align}
with $\epsilon=0.007$, the largest average-RI difference in a two-run pilot. Every selected step
lies in $100$--$200$ (Figure~\ref{fig:ristep}, Table~\ref{tab:selected}), inside the 200-step
warmup and under one pass over the training grid.

\begin{table}[htbp]
\caption{Selected checkpoint step per backbone and adapter seed under the rule above.
Columns are seed identifiers: Midnight-12k, Virchow2 and H-optimus-0 use seeds 0, 1, 3, the
others 0, 1, 2.}
\label{tab:selected}
\centering
\begin{tabular}{lccc}
\toprule
Backbone & Seed 0 & Seed 1 & Seed 2\,/\,3 \\
\midrule
Phikon-v2 & 200 & 200 & 200 \\
Midnight-12k & 150 & 100 & 100 \\
Virchow2 & 100 & 150 & 100 \\
H-optimus-0 & 100 & 100 & 100 \\
UNI2-h & 100 & 150 & 100 \\
Virchow & 150 & 150 & 150 \\
OpenMidnight & 150 & 150 & 150 \\
\bottomrule
\end{tabular}
\end{table}

\begin{figure}[htbp]
\centering
\includegraphics[width=\linewidth]{figures/ri_vs_step.pdf}
\caption{PathoROB RI against step for all 21 runs (three seeds per backbone). Circles mark
the selected checkpoint, the band is steps $100$--$200$, and the dashed line is the base
model. Each panel has its own $y$-range.}
\label{fig:ristep}
\end{figure}

\section{Evaluation protocol details}
\label{app:evaluation}

\paragraph{PathoROB.}
We run the benchmark's own harness unchanged at a pinned revision, so RI follows its published
definition~\citep{koemen2026pathorob} (\S\ref{sec:experiments}): for every query tile the $k$ nearest neighbours are found
by cosine similarity over $\ell_2$-normalized embeddings. Embeddings are the CLS token concatenated with the mean patch token
for every backbone, the pooling of the benchmark's reference rows, without projection heads.
$k=61$, $11$, $46$ for TCGA $2\times2$, Camelyon, Tolkach-ESCA, the harness's per-dataset defaults. \texttt{prediction\_performance}
is the fraction of $k$ nearest neighbours sharing the biological class, pooled over same- and
other-centre neighbours. The submetrics tabulated in Appendix Table~\ref{tab:pathorob-sub}
follow the same harness: balanced accuracy is $k$-NN balanced accuracy at $k_{\mathrm{opt}}$,
and ID, OOD and pooled purity are the fractions of same-centre, other-centre and all $k$
nearest neighbours that share the biological class. All five are higher-is-better. Our base RI reproduces the published rows to $0.004$ or better except
H-optimus-0 (TCGA $0.803$ vs $0.812$, Camelyon $0.679$ vs $0.705$); both H-optimus-0 columns are
our own measurement. Midnight-12k and OpenMidnight are absent from the published field, so they
are ranked by insertion rather than substitution, making that field 24 models rather than 23.

\paragraph{Held-out PLISM retrieval.}
Every condition is sampled at the same 256 coordinates. The embedding concatenates the CLS token and the mean patch token, with no projection head. A pair is two acquisition conditions that differ in scanner alone or in stain alone, and both conditions are held out from training. One supplies the 256 query tiles, the other the 256 candidates at the same coordinates. We rank candidates by cosine similarity after $\ell_2$ normalisation and score top-1 accuracy per pair, then average over pairs. PLISM-wsi spans seven scanners (AT2, GT450, P, S210, S360, S60, SQ) and 13 staining protocols (GIV, GIVH, GM, GMH, GV, GVH, HR, HRH, KR, KRH, LM, LMH, MY, in PLISM's codes). The 41 held-out conditions (\S\ref{sec:data}) comprise 20 with an unseen scanner, 15 with an unseen stain, and six with both. This gives 73 cross-scanner pairs and 171 cross-stain pairs. These aggregates test held-out acquisition conditions; the varied scanner or stain is not necessarily unseen.

\paragraph{HEST.}
We use the official train/test splits and 50-gene list per task. Standardization and
256-dimensional PCA are fit on training embeddings only. Ridge is applied with HEST defaults. We report Pearson $r$ per
gene, averaged over genes within a task, then over nine tasks. Pooling is CLS for Phikon-v2,
Midnight-12k, H-optimus-0, UNI2-h, OpenMidnight, and concatenated CLS and mean-patch for
Virchow2 and Virchow, identical for base and fine-tuned. Base averages reproduce published rows
to $0.0010$.

\paragraph{CPTAC.}
We test models across 38 Patho-Bench \citep{pathobench} CPTAC classification tasks, on the benchmark's published splits: nine on
COAD, five on LUAD, four each on CCRCC and LSCC, three each on BRCA, GBM, HNSC and UCEC, two on
PDA, and one each on LUNG (subtype) and OV. Targets are gene-mutation status, MSI-H, histologic
grade, immune class and subtype. Tile embeddings are mean-pooled to a slide vector and, for
case-level tasks, slides are mean-pooled to a case vector; the probe is the benchmark's
class-balanced logistic regression ($C=0.5$, 100 bootstraps), and the metric is macro
one-vs-rest AUC, averaged without weighting over the 38 tasks (Table~\ref{tab:hestcptac}; per
task in Table~\ref{tab:cptac-per-task}). Pooling is CLS for Phikon-v2, H-optimus-0, UNI2-h, OpenMidnight, and concatenated CLS and mean-patch for
Virchow2 and Virchow, Midnight-12k, identical for base and fine-tuned. The benchmark's four survival tasks are not scored.

\paragraph{THUNDER.}
Both terms of every base $\to$ fine-tuned pair are our own measurement under the benchmark's
harness, and a rank is obtained by substituting that value for the backbone's published row in
the 32-model field, so base and fine-tuned ranks are comparable to each other and to the
published field. The segmentation probe is non-deterministic, so a segmentation spread folds in
harness repeat variability as well as adapter-seed variability, and our measured base can differ
slightly from the published one. ECE is the linear-probe expected calibration error, estimated
from relatively few populated bins, which makes small changes sensitive to how predictions fall
across bins. ECE and the adversarial macro-F1 drop are lower-is-better; every other THUNDER task
is higher-is-better.

\paragraph{Outliers and consistent effects.}
Three patterns dominate the per-dataset panel. TCGA-uniform is the only dataset with a
unanimous macro-F1 regression across two tasks: $k$-NN falls on all
seven backbones ($-10.9$ to $-2.5$, mean $-6.2$, every cell beyond its own two SDs) and linear
probing on all seven ($-6.6$ to $-0.9$, five beyond two SDs). Of the 21 cells spanned by $k$-NN,
linear probing and few-shot, 19 decrease, and those 19 average $4.6$ points. Adversarial
robustness, on the contrary, has the most consistent positive effect in the panel: the F1 drop shrinks
on all seven backbones for 9 of the 16 datasets and worsens unanimously on none, with PCam the
extreme ($-30.8$ to $-12.9$, mean $-19.9$, all seven beyond two SDs) followed by BACH ($-21.1$
to $-7.3$, mean $-11.0$). The ECE degradation is concentrated on the two small breast datasets,
which supply three of the four largest ECE increases exceeding two SDs (the other is Virchow2 on BreakHis, $+7.1$): BRACS worsens on six of seven backbones
($+2.3$ to $+13.5$, three beyond two SDs, the exception being H-optimus-0), and BACH on UNI2-h
moves $8.4\to23.4$, the largest in the panel, while its BACH linear-probe macro-F1 falls $1.8$. Excluding
BRACS, the per-backbone mean $\Delta$ECE over the remaining 15 datasets is $-0.21$, $+0.81$,
$-0.39$, $-0.42$, $+0.44$, $-0.84$, $-1.59$, so only Midnight-12k and UNI2-h still rise.
Segmentation shows no systematic pattern on any of its four datasets. The volatile cells are the
small ones: BACH has 132 test images over four classes with 23 in the smallest, and BRACS 570
over seven, against $32{,}768$ for PCam and $56{,}275$ for TCGA-TILS. TCGA-uniform is large
($56{,}200$) but 32-way, with 180 in its smallest class. Published THUNDER mean ECE spans
$2.9$--$5.8$ across pathology encoders.



\section{Full result tables}
% Full-page benchmark tables: let a single table own a float page even when it fills less than half of it,
% so the tables follow this section instead of being deferred past Appendix~\ref{app:implementation}.
% Left in force for the rest of the appendix: the only later float is the ablations table.
\renewcommand{\floatpagefraction}{0.5}

We report full tables per benchmark, in the order of the main text: PathoROB (Table~\ref{tab:pathorob-sub}; RI rank sums and PLISM retrieval are in main-text Table~\ref{tab:pathorobranks}), HEST (Table~\ref{tab:hestranks}), CPTAC (Table~\ref{tab:cptac-per-task}), and THUNDER (Table~\ref{tab:ranks}, continued in Table~\ref{tab:ranks-b}; Figure~\ref{fig:thunderds}).
Every fine-tuned value in the paper is a mean over three adapter seeds at each seed's selected checkpoint (Table~\ref{tab:selected}). Base values are single evaluations (deterministic except segmentation) and carry no spread. Per-seed summary blocks list the unaggregated values behind the aggregate rows: the base value and each adapter seed's value at its own selected checkpoint, seed columns labelled a--c and ordered by seed identifier within a backbone (0, 1, 2, or 0, 1, 3 for Midnight-12k, Virchow2 and H-optimus-0), so they do not pair the same seed across backbones. Ablations (Appendix~\ref{app:ablations}) are single seed and carry none.
One highlighting rule holds throughout: in every value cell the fine-tuned number is
\gain{bold} when the change from base is an improvement exceeding two sample SDs, \loss{red}
when it is a regression exceeding two sample SDs, and plain when it lies within two SDs, with
improvement read in the metric's own direction (a decrease is an improvement for ECE and for
the adversarial drop). Rank cells have no spread, so a rank is \gain{bold} when it improves and
\loss{red} when it worsens.

On HEST (\S\ref{sec:results}), the largest resolved regression is H-optimus-0 on COAD ($-0.022$ against $\pm0.004$). On THUNDER, the fraction of per-dataset cells whose change exceeds two SDs is $72/112$ for $k$-NN, $59/112$ for linear probing, $94/112$ for few-shot, $8/28$ for segmentation, $55/112$ for ECE and $89/112$ for adversarial drop. Of the six BRACS ECE increases, three (Midnight-12k, Virchow2, Virchow) exceed two SDs and three (Phikon-v2, UNI2-h, OpenMidnight) lie within it.

\begin{table}[p]
\caption{PathoROB. Top: submetrics per dataset, base $\to$ fine-tuned mean $\pm$ two sample
SDs; all five are higher-is-better and are defined in Appendix~\ref{app:evaluation}.
Bottom: per-seed mean RI, base and three adapter seeds.}
\label{tab:pathorob-sub}
\centering
\footnotesize
\setlength{\tabcolsep}{3pt}
\resizebox{\textwidth}{!}{%
\input{tables/pathorob_submetrics}}
\par\medskip
{\footnotesize\emph{Per-seed mean RI}}\par\smallskip
\input{tables/perseed_pathorob}
\end{table}

\begin{table}[p]
\caption{HEST. Top: per task, base $\to$ fine-tuned mean Pearson $r$ $\pm$ two sample SDs on
the upper line and that value's rank in the 26-model published field below it; the Mean $r$ and
$\Sigma$rank columns repeat Table~\ref{tab:hestcptac}.
Bottom: per-seed mean $r$ over the nine tasks.}
\label{tab:hestranks}
\centering
\footnotesize
\setlength{\tabcolsep}{2.5pt}
\resizebox{\textwidth}{!}{%
\input{tables/hest_per_task}}
\par\medskip
{\footnotesize\emph{Per-seed mean $r$}}\par\smallskip
\input{tables/perseed_hest}
\end{table}

\begin{table}[p]
\caption{CPTAC. Top: per task, base $\to$ fine-tuned macro one-vs-rest AUC. Highlighting is
gated on the two-SD spread, which this table does not print; the per-seed block below carries
the values behind it. The Mean row matches Table~\ref{tab:hestcptac}.
Bottom: per-seed macro-OvR AUC.}
\label{tab:cptac-per-task}
\centering
\footnotesize
\setlength{\tabcolsep}{3pt}
\resizebox{\textwidth}{!}{%
\input{tables/cptac_per_task}}
\par\medskip
{\footnotesize\emph{Per-seed macro-OvR AUC}}\par\smallskip
\input{tables/perseed_cptac}
\end{table}

\begin{table}[p]
\caption{THUNDER, part one of two. Top: task means, base $\to$ fine-tuned mean $\pm$ two
sample SDs, each value's rank in the 32-model field in parentheses, then the rank sum and its
implied position; ECE and the adversarial drop are lower-is-better.
Bottom: per-dataset macro-F1 for $k$-NN, linear probing and 16-shot, same spread.
Continued in Table~\ref{tab:ranks-b}.}
\label{tab:ranks}
\centering
\footnotesize
\setlength{\tabcolsep}{3pt}
\resizebox{\textwidth}{!}{%
\input{tables/thunder_ranks_full}}
\par\medskip
{\scriptsize
\setlength{\tabcolsep}{3pt}
\renewcommand{\arraystretch}{0.95}
\resizebox{\textwidth}{!}{%
\input{tables/thunder_per_ds_cls}}}
\end{table}

\begin{table}[p]
\caption{THUNDER, part two of two (continued from Table~\ref{tab:ranks}). Top: per dataset,
segmentation F1, linear-probe ECE and adversarial macro-F1 drop (percentage points, the latter
two lower-is-better; fp32-attack cells), base $\to$ fine-tuned mean $\pm$ two sample SDs; the
segmentation spread also folds in harness repeat variability (Appendix~\ref{app:evaluation}).
Bottom: per-seed task means, percentage points.}
\label{tab:ranks-b}
\centering
{\scriptsize
\setlength{\tabcolsep}{3pt}
\renewcommand{\arraystretch}{0.95}
\resizebox{\textwidth}{!}{%
\input{tables/thunder_per_ds_other}}}
\par\medskip
{\footnotesize\emph{Per-seed task means}}\par\smallskip
\scriptsize
\setlength{\tabcolsep}{3pt}
\resizebox{\textwidth}{!}{%
\input{tables/perseed_thunder}}
\end{table}

\begin{figure}[t]
\centering
\includegraphics[width=\linewidth]{figures/thunder_per_ds.pdf}
\caption{THUNDER. Per-dataset change, fine-tuned $-$ base; bold exceeds two SDs on that
dataset, faint lies within. In the ECE and adversarial panels (lower is better) the colour
scale is reversed, so blue always improves; the adversarial scale is wider. Bottom row of each
panel: mean across backbones. Right-hand columns: row means, the last two dropping
tcga\_uniform and BRACS.}
\label{fig:thunderds}
\end{figure}

\section{Ablation details}
\label{app:ablations}

\paragraph{Pseudo-tissue reweighting.}
Variants: asymmetric $(+3,-\infty)$, symmetric $(+3,+3)$, none $(0,0)$. Symmetric has lower RI
than asymmetric in all 16 step-matched comparisons. At the asymmetric arm's seed-0 selected
checkpoint, asymmetric minus none RI and HEST are in Table~\ref{tab:ablations}(a). THUNDER $k$-NN, linear probing, and few-shot are lower with
asymmetric reweighting for every backbone except Phikon-v2 $k$-NN ($+0.1$). Only Phikon-v2
clears the RI tolerance $\epsilon$, and on HEST only Phikon-v2 and H-optimus-0
exceed the tuned across-seed spread, which is why the recipe does not claim reweighting as a
general mechanism.

\paragraph{Batch composition.}
Phikon-v2, one seed, a diagnostic sweep (RI compared at step 750; fixed-$C$ changes also at steps 1000 and 1250) at a fixed 3,600-image budget with $CT$ held
constant. Table~\ref{tab:ablations}(b) reports RI at step 750 as the number of conditions per
batch $C$ decreases, and, for the two $C$ with a second $T$ run, the change in RI from
increasing the number of locations $T$ at fixed $C$, at three steps. Because $C$ and $T$ covary
at fixed $CT$, the first block is an allocation comparison rather than an isolated effect of
$C$; the recipe uses $C=2$.

\paragraph{Adapter rank.}
Mean RI over checkpoints at the stated steps, at fixed $\alpha/r=2$
(Table~\ref{tab:ablations}(c)). No monotone trend is present on Phikon-v2; on Virchow2 rank 128
is clearly worse, and the recipe uses rank 32. This sweep and the batch-composition sweep both
average over steps beyond the $100$--$200$ range where the recipe's checkpoints are selected
(\S\ref{sec:selection}) and are one-seed diagnostics, so they bound these effects rather than
settle the settings.

\begin{table}[htbp]
\caption{Ablations, one seed each: (a) pseudo-tissue reweighting, (b) batch composition,
(c) adapter rank. Differences are read against the paper's
RI tolerance $\epsilon=0.007$ and, on HEST, against the tuned across-seed $2$\,SD of the same
backbone in the rightmost column of block (a); $^\dagger$ marks $|\Delta\mathrm{RI}|<\epsilon$,
and \textbf{bold} marks the recipe's setting.}
\label{tab:ablations}
\centering
\footnotesize
\setlength{\tabcolsep}{6pt}
\input{tables/ablations}
\end{table}

\clearpage
\section{Implementation details}
\label{app:implementation}

\paragraph{Adapters and heads.}
LoRA rank 32, $\alpha=64$, dropout 0, on every attention projection and MLP layer. Projection
heads are two-layer MLPs. GeM exponent initialized to 3, constrained to $[1,8]$. $\tau=0.07$.
Trainable parameters are $2.7$--$4.4\%$ of the backbone parameters for the LoRA adapters
alone, i.e.\ the parameters actually merged into the released backbone (Phikon-v2 $4.41$,
Midnight-12k $3.01$, Virchow2 and Virchow $3.20$, H-optimus-0 and OpenMidnight $2.69$,
UNI2-h $2.68$), and $3.0$--$5.4\%$ for the adapters plus the two projection heads
(Phikon-v2 $5.40$, Midnight-12k $3.37$, Virchow2 $3.76$, H-optimus-0 $3.05$, UNI2-h $3.28$,
Virchow $3.76$, OpenMidnight $3.05$); the projection heads are discarded before merging.

\paragraph{Optimization and preprocessing.}
Native $224\times224$ tiles, no resizing, cropping, colour, or geometric augmentation; each
backbone's published channel normalization. AdamW, learning rate $10^{-4}$, weight decay $0.05$,
200 linear warmup steps, cosine decay over 500 steps, gradient clipping at $1.0$. bfloat16, gradient checkpointing,
activation offloading. Checkpoints every 50 steps. Three seeds per backbone controlling
initialization and sampling.

\paragraph{Compute.}
Two H100 80GB per job: one trains, one runs concurrent RI evaluation. 500-step runs take
$1.0$--$7.8$ hours per backbone. Evaluation dominates total compute.

\clearpage
\section*{NeurIPS Paper Checklist}

%%% BEGIN INSTRUCTIONS %%%
The checklist is designed to encourage best practices for responsible machine learning research, addressing issues of reproducibility, transparency, research ethics, and societal impact. Do not remove the checklist: {\bf The papers not including the checklist will be desk rejected.} The checklist should follow the references and follow the (optional) supplemental material.  The checklist does NOT count towards the page
limit. 

Please read the checklist guidelines carefully for information on how to answer these questions. For each question in the checklist:
\begin{itemize}
    \item You should answer \answerYes{}, \answerNo{}, or \answerNA{}.
    \item \answerNA{} means either that the question is Not Applicable for that particular paper or the relevant information is Not Available.
    \item Please provide a short (1--2 sentence) justification right after your answer (even for \answerNA). 
   % \item {\bf The papers not including the checklist will be desk rejected.}
\end{itemize}

{\bf The checklist answers are an integral part of your paper submission.} They are visible to the reviewers, area chairs, senior area chairs, and ethics reviewers. You will also be asked to include it (after eventual revisions) with the final version of your paper, and its final version will be published with the paper.

The reviewers of your paper will be asked to use the checklist as one of the factors in their evaluation. While \answerYes{} is generally preferable to \answerNo{}, it is perfectly acceptable to answer \answerNo{} provided a proper justification is given (e.g., error bars are not reported because it would be too computationally expensive'' or ``we were unable to find the license for the dataset we used''). In general, answering \answerNo{} or \answerNA{} is not grounds for rejection. While the questions are phrased in a binary way, we acknowledge that the true answer is often more nuanced, so please just use your best judgment and write a justification to elaborate. All supporting evidence can appear either in the main paper or the supplemental material, provided in appendix. If you answer \answerYes{} to a question, in the justification please point to the section(s) where related material for the question can be found.

IMPORTANT, please:
\begin{itemize}
    \item {\bf Delete this instruction block, but keep the section heading ``NeurIPS Paper Checklist"},
    \item  {\bf Keep the checklist subsection headings, questions/answers and guidelines below.}
    \item {\bf Do not modify the questions and only use the provided macros for your answers}.
\end{itemize} 
 

%%% END INSTRUCTIONS %%%


\begin{enumerate}

\item {\bf Claims}
    \item[] Question: Do the main claims made in the abstract and introduction accurately reflect the paper's contributions and scope?
    \item[] Answer: \answerTODO{} % Replace by \answerYes{}, \answerNo{}, or \answerNA{}.
    \item[] Justification: \justificationTODO{}
    \item[] Guidelines:
    \begin{itemize}
        \item The answer \answerNA{} means that the abstract and introduction do not include the claims made in the paper.
        \item The abstract and/or introduction should clearly state the claims made, including the contributions made in the paper and important assumptions and limitations. A \answerNo{} or \answerNA{} answer to this question will not be perceived well by the reviewers. 
        \item The claims made should match theoretical and experimental results, and reflect how much the results can be expected to generalize to other settings. 
        \item It is fine to include aspirational goals as motivation as long as it is clear that these goals are not attained by the paper. 
    \end{itemize}

\item {\bf Limitations}
    \item[] Question: Does the paper discuss the limitations of the work performed by the authors?
    \item[] Answer: \answerTODO{} % Replace by \answerYes{}, \answerNo{}, or \answerNA{}.
    \item[] Justification: \justificationTODO{}
    \item[] Guidelines:
    \begin{itemize}
        \item The answer \answerNA{} means that the paper has no limitation while the answer \answerNo{} means that the paper has limitations, but those are not discussed in the paper. 
        \item The authors are encouraged to create a separate ``Limitations'' section in their paper.
        \item The paper should point out any strong assumptions and how robust the results are to violations of these assumptions (e.g., independence assumptions, noiseless settings, model well-specification, asymptotic approximations only holding locally). The authors should reflect on how these assumptions might be violated in practice and what the implications would be.
        \item The authors should reflect on the scope of the claims made, e.g., if the approach was only tested on a few datasets or with a few runs. In general, empirical results often depend on implicit assumptions, which should be articulated.
        \item The authors should reflect on the factors that influence the performance of the approach. For example, a facial recognition algorithm may perform poorly when image resolution is low or images are taken in low lighting. Or a speech-to-text system might not be used reliably to provide closed captions for online lectures because it fails to handle technical jargon.
        \item The authors should discuss the computational efficiency of the proposed algorithms and how they scale with dataset size.
        \item If applicable, the authors should discuss possible limitations of their approach to address problems of privacy and fairness.
        \item While the authors might fear that complete honesty about limitations might be used by reviewers as grounds for rejection, a worse outcome might be that reviewers discover limitations that aren't acknowledged in the paper. The authors should use their best judgment and recognize that individual actions in favor of transparency play an important role in developing norms that preserve the integrity of the community. Reviewers will be specifically instructed to not penalize honesty concerning limitations.
    \end{itemize}

\item {\bf Theory assumptions and proofs}
    \item[] Question: For each theoretical result, does the paper provide the full set of assumptions and a complete (and correct) proof?
    \item[] Answer: \answerTODO{} % Replace by \answerYes{}, \answerNo{}, or \answerNA{}.
    \item[] Justification: \justificationTODO{}
    \item[] Guidelines:
    \begin{itemize}
        \item The answer \answerNA{} means that the paper does not include theoretical results. 
        \item All the theorems, formulas, and proofs in the paper should be numbered and cross-referenced.
        \item All assumptions should be clearly stated or referenced in the statement of any theorems.
        \item The proofs can either appear in the main paper or the supplemental material, but if they appear in the supplemental material, the authors are encouraged to provide a short proof sketch to provide intuition. 
        \item Inversely, any informal proof provided in the core of the paper should be complemented by formal proofs provided in appendix or supplemental material.
        \item Theorems and Lemmas that the proof relies upon should be properly referenced. 
    \end{itemize}

    \item {\bf Experimental result reproducibility}
    \item[] Question: Does the paper fully disclose all the information needed to reproduce the main experimental results of the paper to the extent that it affects the main claims and/or conclusions of the paper (regardless of whether the code and data are provided or not)?
    \item[] Answer: \answerTODO{} % Replace by \answerYes{}, \answerNo{}, or \answerNA{}.
    \item[] Justification: \justificationTODO{}
    \item[] Guidelines:
    \begin{itemize}
        \item The answer \answerNA{} means that the paper does not include experiments.
        \item If the paper includes experiments, a \answerNo{} answer to this question will not be perceived well by the reviewers: Making the paper reproducible is important, regardless of whether the code and data are provided or not.
        \item If the contribution is a dataset and\slash or model, the authors should describe the steps taken to make their results reproducible or verifiable. 
        \item Depending on the contribution, reproducibility can be accomplished in various ways. For example, if the contribution is a novel architecture, describing the architecture fully might suffice, or if the contribution is a specific model and empirical evaluation, it may be necessary to either make it possible for others to replicate the model with the same dataset, or provide access to the model. In general. releasing code and data is often one good way to accomplish this, but reproducibility can also be provided via detailed instructions for how to replicate the results, access to a hosted model (e.g., in the case of a large language model), releasing of a model checkpoint, or other means that are appropriate to the research performed.
        \item While NeurIPS does not require releasing code, the conference does require all submissions to provide some reasonable avenue for reproducibility, which may depend on the nature of the contribution. For example
        \begin{enumerate}
            \item If the contribution is primarily a new algorithm, the paper should make it clear how to reproduce that algorithm.
            \item If the contribution is primarily a new model architecture, the paper should describe the architecture clearly and fully.
            \item If the contribution is a new model (e.g., a large language model), then there should either be a way to access this model for reproducing the results or a way to reproduce the model (e.g., with an open-source dataset or instructions for how to construct the dataset).
            \item We recognize that reproducibility may be tricky in some cases, in which case authors are welcome to describe the particular way they provide for reproducibility. In the case of closed-source models, it may be that access to the model is limited in some way (e.g., to registered users), but it should be possible for other researchers to have some path to reproducing or verifying the results.
        \end{enumerate}
    \end{itemize}


\item {\bf Open access to data and code}
    \item[] Question: Does the paper provide open access to the data and code, with sufficient instructions to faithfully reproduce the main experimental results, as described in supplemental material?
    \item[] Answer: \answerTODO{} % Replace by \answerYes{}, \answerNo{}, or \answerNA{}.
    \item[] Justification: \justificationTODO{}
    \item[] Guidelines:
    \begin{itemize}
        \item The answer \answerNA{} means that paper does not include experiments requiring code.
        \item Please see the NeurIPS code and data submission guidelines (\url{https://neurips.cc/public/guides/CodeSubmissionPolicy}) for more details.
        \item While we encourage the release of code and data, we understand that this might not be possible, so \answerNo{} is an acceptable answer. Papers cannot be rejected simply for not including code, unless this is central to the contribution (e.g., for a new open-source benchmark).
        \item The instructions should contain the exact command and environment needed to run to reproduce the results. See the NeurIPS code and data submission guidelines (\url{https://neurips.cc/public/guides/CodeSubmissionPolicy}) for more details.
        \item The authors should provide instructions on data access and preparation, including how to access the raw data, preprocessed data, intermediate data, and generated data, etc.
        \item The authors should provide scripts to reproduce all experimental results for the new proposed method and baselines. If only a subset of experiments are reproducible, they should state which ones are omitted from the script and why.
        \item At submission time, to preserve anonymity, the authors should release anonymized versions (if applicable).
        \item Providing as much information as possible in supplemental material (appended to the paper) is recommended, but including URLs to data and code is permitted.
    \end{itemize}


\item {\bf Experimental setting/details}
    \item[] Question: Does the paper specify all the training and test details (e.g., data splits, hyperparameters, how they were chosen, type of optimizer) necessary to understand the results?
    \item[] Answer: \answerTODO{} % Replace by \answerYes{}, \answerNo{}, or \answerNA{}.
    \item[] Justification: \justificationTODO{}
    \item[] Guidelines:
    \begin{itemize}
        \item The answer \answerNA{} means that the paper does not include experiments.
        \item The experimental setting should be presented in the core of the paper to a level of detail that is necessary to appreciate the results and make sense of them.
        \item The full details can be provided either with the code, in appendix, or as supplemental material.
    \end{itemize}

\item {\bf Experiment statistical significance}
    \item[] Question: Does the paper report error bars suitably and correctly defined or other appropriate information about the statistical significance of the experiments?
    \item[] Answer: \answerTODO{} % Replace by \answerYes{}, \answerNo{}, or \answerNA{}.
    \item[] Justification: \justificationTODO{}
    \item[] Guidelines:
    \begin{itemize}
        \item The answer \answerNA{} means that the paper does not include experiments.
        \item The authors should answer \answerYes{} if the results are accompanied by error bars, confidence intervals, or statistical significance tests, at least for the experiments that support the main claims of the paper.
        \item The factors of variability that the error bars are capturing should be clearly stated (for example, train/test split, initialization, random drawing of some parameter, or overall run with given experimental conditions).
        \item The method for calculating the error bars should be explained (closed form formula, call to a library function, bootstrap, etc.)
        \item The assumptions made should be given (e.g., Normally distributed errors).
        \item It should be clear whether the error bar is the standard deviation or the standard error of the mean.
        \item It is OK to report 1-sigma error bars, but one should state it. The authors should preferably report a 2-sigma error bar than state that they have a 96\% CI, if the hypothesis of Normality of errors is not verified.
        \item For asymmetric distributions, the authors should be careful not to show in tables or figures symmetric error bars that would yield results that are out of range (e.g., negative error rates).
        \item If error bars are reported in tables or plots, the authors should explain in the text how they were calculated and reference the corresponding figures or tables in the text.
    \end{itemize}

\item {\bf Experiments compute resources}
    \item[] Question: For each experiment, does the paper provide sufficient information on the computer resources (type of compute workers, memory, time of execution) needed to reproduce the experiments?
    \item[] Answer: \answerTODO{} % Replace by \answerYes{}, \answerNo{}, or \answerNA{}.
    \item[] Justification: \justificationTODO{}
    \item[] Guidelines:
    \begin{itemize}
        \item The answer \answerNA{} means that the paper does not include experiments.
        \item The paper should indicate the type of compute workers CPU or GPU, internal cluster, or cloud provider, including relevant memory and storage.
        \item The paper should provide the amount of compute required for each of the individual experimental runs as well as estimate the total compute. 
        \item The paper should disclose whether the full research project required more compute than the experiments reported in the paper (e.g., preliminary or failed experiments that didn't make it into the paper). 
    \end{itemize}
    
\item {\bf Code of ethics}
    \item[] Question: Does the research conducted in the paper conform, in every respect, with the NeurIPS Code of Ethics \url{https://neurips.cc/public/EthicsGuidelines}?
    \item[] Answer: \answerTODO{} % Replace by \answerYes{}, \answerNo{}, or \answerNA{}.
    \item[] Justification: \justificationTODO{}
    \item[] Guidelines:
    \begin{itemize}
        \item The answer \answerNA{} means that the authors have not reviewed the NeurIPS Code of Ethics.
        \item If the authors answer \answerNo, they should explain the special circumstances that require a deviation from the Code of Ethics.
        \item The authors should make sure to preserve anonymity (e.g., if there is a special consideration due to laws or regulations in their jurisdiction).
    \end{itemize}


\item {\bf Broader impacts}
    \item[] Question: Does the paper discuss both potential positive societal impacts and negative societal impacts of the work performed?
    \item[] Answer: \answerTODO{} % Replace by \answerYes{}, \answerNo{}, or \answerNA{}.
    \item[] Justification: \justificationTODO{}
    \item[] Guidelines:
    \begin{itemize}
        \item The answer \answerNA{} means that there is no societal impact of the work performed.
        \item If the authors answer \answerNA{} or \answerNo, they should explain why their work has no societal impact or why the paper does not address societal impact.
        \item Examples of negative societal impacts include potential malicious or unintended uses (e.g., disinformation, generating fake profiles, surveillance), fairness considerations (e.g., deployment of technologies that could make decisions that unfairly impact specific groups), privacy considerations, and security considerations.
        \item The conference expects that many papers will be foundational research and not tied to particular applications, let alone deployments. However, if there is a direct path to any negative applications, the authors should point it out. For example, it is legitimate to point out that an improvement in the quality of generative models could be used to generate Deepfakes for disinformation. On the other hand, it is not needed to point out that a generic algorithm for optimizing neural networks could enable people to train models that generate Deepfakes faster.
        \item The authors should consider possible harms that could arise when the technology is being used as intended and functioning correctly, harms that could arise when the technology is being used as intended but gives incorrect results, and harms following from (intentional or unintentional) misuse of the technology.
        \item If there are negative societal impacts, the authors could also discuss possible mitigation strategies (e.g., gated release of models, providing defenses in addition to attacks, mechanisms for monitoring misuse, mechanisms to monitor how a system learns from feedback over time, improving the efficiency and accessibility of ML).
    \end{itemize}
    
\item {\bf Safeguards}
    \item[] Question: Does the paper describe safeguards that have been put in place for responsible release of data or models that have a high risk for misuse (e.g., pre-trained language models, image generators, or scraped datasets)?
    \item[] Answer: \answerTODO{} % Replace by \answerYes{}, \answerNo{}, or \answerNA{}.
    \item[] Justification: \justificationTODO{}
    \item[] Guidelines:
    \begin{itemize}
        \item The answer \answerNA{} means that the paper poses no such risks.
        \item Released models that have a high risk for misuse or dual-use should be released with necessary safeguards to allow for controlled use of the model, for example by requiring that users adhere to usage guidelines or restrictions to access the model or implementing safety filters. 
        \item Datasets that have been scraped from the Internet could pose safety risks. The authors should describe how they avoided releasing unsafe images.
        \item We recognize that providing effective safeguards is challenging, and many papers do not require this, but we encourage authors to take this into account and make a best faith effort.
    \end{itemize}

\item {\bf Licenses for existing assets}
    \item[] Question: Are the creators or original owners of assets (e.g., code, data, models), used in the paper, properly credited and are the license and terms of use explicitly mentioned and properly respected?
    \item[] Answer: \answerTODO{} % Replace by \answerYes{}, \answerNo{}, or \answerNA{}.
    \item[] Justification: \justificationTODO{}
    \item[] Guidelines:
    \begin{itemize}
        \item The answer \answerNA{} means that the paper does not use existing assets.
        \item The authors should cite the original paper that produced the code package or dataset.
        \item The authors should state which version of the asset is used and, if possible, include a URL.
        \item The name of the license (e.g., CC-BY 4.0) should be included for each asset.
        \item For scraped data from a particular source (e.g., website), the copyright and terms of service of that source should be provided.
        \item If assets are released, the license, copyright information, and terms of use in the package should be provided. For popular datasets, \url{paperswithcode.com/datasets} has curated licenses for some datasets. Their licensing guide can help determine the license of a dataset.
        \item For existing datasets that are re-packaged, both the original license and the license of the derived asset (if it has changed) should be provided.
        \item If this information is not available online, the authors are encouraged to reach out to the asset's creators.
    \end{itemize}

\item {\bf New assets}
    \item[] Question: Are new assets introduced in the paper well documented and is the documentation provided alongside the assets?
    \item[] Answer: \answerTODO{} % Replace by \answerYes{}, \answerNo{}, or \answerNA{}.
    \item[] Justification: \justificationTODO{}
    \item[] Guidelines:
    \begin{itemize}
        \item The answer \answerNA{} means that the paper does not release new assets.
        \item Researchers should communicate the details of the dataset\slash code\slash model as part of their submissions via structured templates. This includes details about training, license, limitations, etc. 
        \item The paper should discuss whether and how consent was obtained from people whose asset is used.
        \item At submission time, remember to anonymize your assets (if applicable). You can either create an anonymized URL or include an anonymized zip file.
    \end{itemize}

\item {\bf Crowdsourcing and research with human subjects}
    \item[] Question: For crowdsourcing experiments and research with human subjects, does the paper include the full text of instructions given to participants and screenshots, if applicable, as well as details about compensation (if any)? 
    \item[] Answer: \answerTODO{} % Replace by \answerYes{}, \answerNo{}, or \answerNA{}.
    \item[] Justification: \justificationTODO{}
    \item[] Guidelines:
    \begin{itemize}
        \item The answer \answerNA{} means that the paper does not involve crowdsourcing nor research with human subjects.
        \item Including this information in the supplemental material is fine, but if the main contribution of the paper involves human subjects, then as much detail as possible should be included in the main paper. 
        \item According to the NeurIPS Code of Ethics, workers involved in data collection, curation, or other labor should be paid at least the minimum wage in the country of the data collector. 
    \end{itemize}

\item {\bf Institutional review board (IRB) approvals or equivalent for research with human subjects}
    \item[] Question: Does the paper describe potential risks incurred by study participants, whether such risks were disclosed to the subjects, and whether Institutional Review Board (IRB) approvals (or an equivalent approval/review based on the requirements of your country or institution) were obtained?
    \item[] Answer: \answerTODO{} % Replace by \answerYes{}, \answerNo{}, or \answerNA{}.
    \item[] Justification: \justificationTODO{}
    \item[] Guidelines:
    \begin{itemize}
        \item The answer \answerNA{} means that the paper does not involve crowdsourcing nor research with human subjects.
        \item Depending on the country in which research is conducted, IRB approval (or equivalent) may be required for any human subjects research. If you obtained IRB approval, you should clearly state this in the paper. 
        \item We recognize that the procedures for this may vary significantly between institutions and locations, and we expect authors to adhere to the NeurIPS Code of Ethics and the guidelines for their institution. 
        \item For initial submissions, do not include any information that would break anonymity (if applicable), such as the institution conducting the review.
    \end{itemize}

\item {\bf Declaration of LLM usage}
    \item[] Question: Does the paper describe the usage of LLMs if it is an important, original, or non-standard component of the core methods in this research? Note that if the LLM is used only for writing, editing, or formatting purposes and does \emph{not} impact the core methodology, scientific rigor, or originality of the research, declaration is not required.
    %this research? 
    \item[] Answer: \answerTODO{} % Replace by \answerYes{}, \answerNo{}, or \answerNA{}.
    \item[] Justification: \justificationTODO{}
    \item[] Guidelines:
    \begin{itemize}
        \item The answer \answerNA{} means that the core method development in this research does not involve LLMs as any important, original, or non-standard components.
        \item Please refer to our LLM policy in the NeurIPS handbook for what should or should not be described.
    \end{itemize}

\end{enumerate}

\end{document}
